\item[\textbf{Problem 14-6}]
Professor Blutarsky is consulting for the president of a corporation that is
planning a company party. The company has a hierarchical structure, that is, the
supervisor relation forms a tree rooted at the president. The human resources
department has ranked each employee with a conviviality rating, which is a real
number. In order to make the party fun for all attendees, the president does not
want both an employee and his or her immediate supervisor to attend.

Professor Blutarsky is given the tree that describes the structure of the
corporation, using the left-child, right-sibling representation described in
Section 10.3. Each node of the tree holds, in addition to the pointers, the name
of an employee and that employee's conviviality ranking. Describe an algorithm
to make up a guest list that maximizes the sum of the conviviality ratings of
the guests. Analyze the running time of your algorithm. 

\Solution

The solution to this problem utilizes Dynamic Programming on a tree.

\subsubsection*{Recurrence Relation}
For any given node, we have two choices to maximize the sum of the conviviality
ratings:
\begin{enumerate}
    \item If the node is included, we cannot include its immediate children. We
    must sum the node's conviviality with the optimal solutions for its
    grandchildren.

    \item If the node is NOT included, we sum the optimal solutions for its
    immediate children.
\end{enumerate}

This can be mathematically expressed as:
$$
    C[\text{node}] = \max
    \begin{pmatrix} 
        \displaystyle \sum_{y \text{ is a child of } \text{node}} C[y]                                 \\[2em] 
        \displaystyle \text{node.conviviality} + \sum_{y \text{ is a grandchild of } \text{node}} C[y] 
    \end{pmatrix}
$$

\subsubsection*{Algorithm Description}

To implement this efficiently using the provided \textbf{left-child,
right-sibling} representation, we perform a recursive \textbf{post-order
traversal} of the tree. This guarantees that before we calculate
$C[\text{node}]$, the optimal values $C[y]$ for all its children and
grandchildren have already been computed. 

We use \textbf{memoization} by storing the computed maximum conviviality value
in a new field, \texttt{C}, on each node so it is only calculated once.

For each visited \texttt{node} in the post-order traversal:
\begin{itemize}
    \item \textbf{Summing Children:} Follow the \texttt{node.left\_child}
    pointer, then iterate through the \texttt{right\_sibling} pointers until
    \texttt{null} is reached. Sum the stored \texttt{C} values of these nodes.

    \item \textbf{Summing Grandchildren:} While iterating through the children
    above, follow each child's \texttt{left\_child} pointer and iterate through
    their respective \texttt{right\_sibling} pointers. Sum the stored \texttt{C}
    values of these nodes.

    \item Apply the recurrence relation to find the maximum of the two choices.
    
    \item Store this maximum in \texttt{node.C}.
\end{itemize}
The final maximum conviviality rating for the entire company party will be
stored in \texttt{root.C}.


\subsubsection*{Time Complexity}
The time complexity of this procedure is $\mathbf{\Theta(n)}$.

\begin{itemize}
    \item Because we use memoization with a post-order traversal, the subproblem
    for each node is solved exactly once.

    \item At each node, finding the children and grandchildren involves
    traversing the \texttt{left\_child} and \texttt{right\_sibling} pointers. 

    \item Across the entire algorithm, every pointer in the tree is followed a
    constant number of times (once when treating it as a child, and once when
    treating it as a grandchild). 

    \item Since a tree with $n$ nodes contains exactly $n-1$ edges (pointers),
    the total amount of work done is directly proportional to the number of
    nodes. Therefore, the time complexity is $\Theta(n)$.
\end{itemize}
